% Faithful typesetting transcription of AppendixA, formal pp23:36--23:37.
\section*{AppendixA: Strongly explicit high-girth graphs}
In this section, we will examine the construction of bipartite high-girth graphs of Lazebnik and Ustimenko [27], and verify that it is indeed strongly explicit, thereby proving Theorem2.12.
Let $q$ be an odd prime power. Let $P$ and $L$ be two infinite sequences of elements from $\mathbb F_q$ indexed as follows:
\begin{align*}
P&=p_1,p_{1,1},p_{1,2},p_{2,1},p_{2,2},p'_{2,2},p_{3,2},\ldots,p_{i,i+1},p_{i+1,i},p_{i+1,i+1},p'_{i+1,i+1},p_{i+1,i+2},\ldots,\\
L&=l_1,l_{1,1},l_{1,2},l_{2,1},l_{2,2},l'_{2,2},l_{3,2},\ldots,l_{i,i+1},l_{i+1,i},l_{i+1,i+1},l'_{i+1,i+1},l_{i+1,i+2},\ldots.
\end{align*}
The names $P$ and $L$ mean points and lines, respectively, due to certain geometric intuition of the construction. We say $P$ is incident with $L$ if they satisfy the following set of equations
\begin{equation}\tag{3}
E_1=\begin{cases}l_{1,1}-p_{1,1}=l_1p_1\\l_{1,2}-p_{1,2}=l_{1,1}p_1\\l_{2,1}-p_{2,1}=l_1p_{1,1},\end{cases}\qquad
E_i=\begin{cases}l_{i,i}-p_{i,i}=l_1p_{i-1,i}\\l'_{i,i}-p'_{i,i}=l_{i,i-1}p_1\\l_{i,i+1}-p_{i,i+1}=l_{i,i}p_1\\l_{i+1,i}-p_{i+1,i}=l_1p'_{i,i}\end{cases}\quad(\forall i\ge2).
\end{equation}
One can verify that for every point $P$, a line $L$ incident with it is uniquely determined by $l_1$, since all other coordinates of $L$ can be computed from the equations above iteratively. Similarly for every line $L$, a point $P$ incident with it is uniquely determined by $p_1$.
Let $P_k\in\mathbb F_q^k$ be the length-$k$ prefix of $P$, and $L_k\in\mathbb F_q^k$ be the length-$k$ prefix of $L$. We define $D(k,q)=(V_1,V_2,E\subseteq V_1\times V_2)$ to be the following bipartite graph:
\begin{itemize}
\item $V_1=V_2=\mathbb F_q^k$;
\item $(u,v)\in E\Longleftrightarrow\exists(P,L)$, $P$ is incident with $L$, $u=P_k$, and $v=L_k$.
\end{itemize}
Equivalently, $u$ and $v$ are connected if and only if they satisfy the first $k$ equations of $\bigcup_i E_i$ in (3). Note that $|V_1|=|V_2|=q^k$, $|E|=q^{k+1}$, and the graph is $q$-regular.
\begin{authorclaim}{TheoremA.1 (Lazebnik and Ustimenko [27])}
Let $k\ge3$ be an odd integer and $q$ be an odd prime power. Then the girth of $D(k,q)$ is at least $k+5$.
\end{authorclaim}
For our purpose, it is sufficient to consider $q=p^r$ for prime $p=O(1)$. To construct the graph explicitly, we need to evaluation field operations over $\mathbb F_q$ for a prime power $q$ in $\operatorname{polylog}(q)$ time, for which we need a representation of $\mathbb F_q$. Recall that $\mathbb F_q$ is isomorphic to $\mathbb F_p[x]/(Q(x))$ for any irreducible degree-$r$ polynomial $Q\in\mathbb F_p[x]$. So the explicit representation of $\mathbb F_q$ follows from the construction of irreducible polynomials by Shoup [39].
\begin{authorclaim}{TheoremA.2 (Shoup [39])}
There is a deterministic algorithm that constructs a degree-$n$ irreducible polynomial over $\mathbb F_p$ given an integer $n$ and a prime power $p$ in $\operatorname{poly}(n,p)$ time.
\end{authorclaim}
Now we check that $D(k,q)$ is strongly explicit, in the sense that given the index $i$ of an edge, we can obtain the indices $j_1$ and $j_2$ in $\operatorname{poly}(k,\log q)$ time such that the $i$th edge connects the $j_1$th vertex in $V_1$ and the $j_2$th vertex in $V_2$. Let $q=p^r$ for a prime $p=O(1)$. We number the vertices and edges as follows.
\begin{enumerate}
\item We identify the elements in the finite field $\mathbb F_q$ as length-$r$ vectors from $[p]^r$, and number all the elements in the lexicographic order.
\item We identify the vertices in both $V_1$ and $V_2$ (that are length $k$ sequences of elements in $\mathbb F_q$) as length-$kr$ vectors from $[p]^{kr}$, and number them in the lexicographic order.
\item For any $i\in[q^k]$ and $j\in[q]$, the $(iq+j)$th edge connects the $i$th vertex $P$ in $V_1$ and the unique vertex $L=(l_1,\ldots)$ in $V_2$ connected to $P$ such that $l_1$ is the $j$th element in $\mathbb F_q$.
\end{enumerate}
Under such a numbering scheme, given the index $iq+j$ of an edge with $i\in[q^k]$ and $j\in[q]$, we can easily determine its two endpoints in $\operatorname{poly}(k,\log q)$ time. Now we are ready to prove Theorem2.12.
\begin{authorclaim}{Remainder of Theorem2.12}
Let $r=r(n)=n^{o(1)}$ be a parameter. For every sufficiently large $n$, there exists an $m=\Theta(n/r)$ and a regular graph $G_{m,n}$ with $m$ vertices, $n$ edges, and girth $\Omega(\log n/\log r)$. Moreover, there exists a $\operatorname{polylog}(n)$-time algorithm $A(n,i)$ for $i\in[n]$ that outputs the indices of the two endpoints of the $i$th edge in $G_{m,n}$, and a $\operatorname{polylog}(n)$-time algorithm $B(n,i,j)$ for $i\in[m]$ that outputs the $j$th edge attaching to the $i$th vertex.
\end{authorclaim}
\begin{proof}
Let $q$ be the power of 3 in the interval $[r,3r)$. Applying TheoremA.2, we construct a fixed representation of $\mathbb F_q$ in $\operatorname{polylog}(q)$ time so that the field operations over $\mathbb F_q$ can be evaluated in $\operatorname{polylog}(q)$ time. Let $k\triangleq\lfloor\log_qn\rfloor-1$ (i.e., $k$ is the largest integer such that $q^{k+1}\le n$), $\ell\triangleq\lceil n/q^{k+1}\rceil$, and $m\triangleq2q^k\ell$. Note that
\begin{align*}
m&\le2q^k\left(\frac n{q^{k+1}}+1\right)=\frac{2n}q+\frac{2q^{k+1}}q\le O(n/r),\\
m&\ge2q^k\cdot\frac n{q^{k+1}}\ge\Omega(n/r),
\end{align*}
and therefore $m=\Theta(n/r)$. Now we construct a graph $G_{m,n}$ with $\ell$ connected components as follows. Each of the first $\ell-1$ connected components is a copy of $D(q,k)$ with $2q^k$ vertices and $q^{k+1}$ edges. The last connected component is a subgraph of $D(q,k)$ with $2q^k$ vertices but only the first $n-q^{k+1}(\ell-1)$ edges.
By TheoremA.1, it is easy to see that the girth of $G_{m,n}$ is at least $k=\Omega(\log n/\log q)=\Omega(\log n/\log r)$. According to our numbering scheme and related discussions, it is also easy to verify that given an edge index $i$, we can compute the two endpoints of the $i$th edge in $G_{m,n}$ in $\operatorname{polylog}(n)$ time. Furthermore, the $j$th edge attaching to the $i$th vertex can be easily computed since each of the $q$ edges attaching to the $i$th vertex is uniquely determined by $l_1\in\mathbb F_q$ or $p_1\in\mathbb F_q$ (see Equation3).
\end{proof}
